% theme: scale_drawings_and_angle_bisectors
\MajorQuestionLesson{縮図と角の二等分線の定理}
\SetRowSpace{1mm}


\TypeQuestion{}{similarity_application_map_scale}
\FigProblemBlock{27mm}{fig/lesson_06_map_scale}{
\Text{縮尺 $1:25000$ の地図上で、地点 $\mathrm{A}$ から地点 $\mathrm{B}$ を通って地点 $\mathrm{C}$ まで進む道を考える。地図上の $\Lseg{AB}$ の長さは $3.6\mathrm{cm}$、$\Lseg{BC}$ の長さは $2.4\mathrm{cm}$ である。ただし、図は略図である。}
\QQ{実際の $\Lseg{AB}$ の道のりを $\mathrm{m}$ で求めなさい。}{$900\mathrm{m}$}
\QQ{地点 $\mathrm{A}$ から地点 $\mathrm{B}$ を通って地点 $\mathrm{C}$ まで進む実際の道のりを $\mathrm{km}$ で求めなさい。}{$1.5\mathrm{km}$}
}{
\QQ{同じ道を縮尺 $1:10000$ の地図に表すと、$\Lseg{AB}$ と $\Lseg{BC}$ を合わせた地図上の長さは何 $\mathrm{cm}$ になるか求めなさい。}{$15\mathrm{cm}$}
}

\TypeQuestion{}{similarity_application_river_scale_drawing}
\FigProblemBlock{45mm}{fig/lesson_06_river_scale}{
\Text{川岸の2点 $\mathrm{A},\mathrm{B}$ と、対岸の点 $\mathrm{P}$ の位置を調べ、縮尺 $1:400$ の縮図をかいた。点 $\mathrm{P}$ から直線 $\Lseg{AB}$ に引いた垂線の足を $\mathrm{H}$ とする。縮図上の $\Lseg{AB}$ の長さは $4.5\mathrm{cm}$、$\Lseg{PH}$ の長さは $2.8\mathrm{cm}$ である。川の両岸は平行で、川幅は一定とする。}
\QQ{実際の川幅 $\Lseg{PH}$ を $\mathrm{m}$ で求めなさい。}{$11.2\mathrm{m}$}
\QQ{実際の $\Lseg{AB}$ の長さを $\mathrm{m}$ で求めなさい。}{$18\mathrm{m}$}
}{
\QQ{同じ場所を縮尺 $1:1000$ の縮図に表すと、線分 $\Lseg{PH}$ の長さは何 $\mathrm{cm}$ になるか求めなさい。}{$1.12\mathrm{cm}$}
}

\TypeQuestion{}{similarity_application_building_scale_drawing}
\FigProblemBlock{45mm}{fig/lesson_06_building_scale}{
\Text{水平な地面に建つ建物の高さを調べる。建物の根元を $\mathrm{B}$、頂点を $\mathrm{T}$ とし、地面上の点 $\mathrm{A}$ に立った人の目の位置を $\mathrm{E}$ とする。目の高さは地面から $1.5\mathrm{m}$ である。線分 $\Lseg{BT}$ 上に、点 $\mathrm{E}$ と同じ高さの点 $\mathrm{H}$ をとった。$\Lseg{EH}$ は地面と平行であり、建物は地面に垂直である。この位置関係を縮尺 $1:200$ の縮図に表すと、$\Lseg{EH}=7.5\mathrm{cm}$、$\Lseg{HT}=4.8\mathrm{cm}$ であった。}
\QQ{実際の $\Lseg{HT}$ の長さを $\mathrm{m}$ で求めなさい。}{$9.6\mathrm{m}$}
}{
\QQ{人が立っている地点から建物の根元までの水平な距離 $\Lseg{AB}$ を $\mathrm{m}$ で求めなさい。}{$15\mathrm{m}$}
\QQ{建物の高さ $\Lseg{BT}$ を $\mathrm{m}$ で求めなさい。}{$11.1\mathrm{m}$}
}

\AnswerColumnBreak
\LessonQuestionBreak

\TypeQuestion{$\Tri{ABC}$ で、$\Lseg{AB}=10\mathrm{cm}$、$\Lseg{AC}=15\mathrm{cm}$、$\Lseg{BC}=20\mathrm{cm}$ とする。$\Ang{BAC}$ の二等分線と辺 $\Lseg{BC}$ の交点を $\mathrm{D}$、$\Ang{ABC}$ の二等分線と線分 $\Lseg{AD}$ の交点を $\mathrm{I}$ とする。次の問いに答えなさい。}{angle_bisector_theorem_two_bisectors}
% 内心・円の性質は使わず、ABCとABDのそれぞれで角の二等分線の定理を使う。
\FigProblemBlock{50mm}{fig/lesson_06_two_bisectors}{
\QQ{$\Lseg{BD},\Lseg{DC}$ の長さをそれぞれ求めなさい。}{$\Lseg{BD}=8\mathrm{cm}$、\par $\Lseg{DC}=12\mathrm{cm}$}
\QQ{$\Lseg{AI}:\Lseg{ID}$ を求めなさい。}{$5:4$}
}{
}

% 長さは図中に集約する。二等分する頂点・図の向き・未知の置き方を混在させる。
\TypeQuestion{次の図で、同じ印をつけた角は等しいものとする。$x$ の値を求めなさい。ただし、長さの単位は $\mathrm{cm}$ とする。}{angle_bisector_theorem_calculation}
\QQFigRow[46mm]{fig/lesson_06_bisector_calc_01}{}{$9$}{fig/lesson_06_bisector_calc_07}{}{$2$}
\QQFigRow[46mm]{fig/lesson_06_bisector_calc_04}{}{$10$}{fig/lesson_06_bisector_calc_09}{}{$3$}


